We applied our proposed method to a cohort of patients treated for CLTI in a single academic medical center. The dataset was derived  from the Carolina Data Warehouse for Health, which is housed within the North Carolina Translational and Clinical Sciences Institute (NC TraCS) and includes electronic health record data derived from care provided at the University of North Carolina at Chapel Hill Medical Center. The cohort included patients that met the hemodynamic and symptomatic diagnostic criteria for CLTI and underwent an open surgical or endovascular revascularization to restore blood flow to the lower limbs between April 2013 and October 2017. Details for this cohort were described previously \citep{mcginigle_precision_2021}.

CLTI predominantly affects older adults, who have a history of smoking or other risk factors for developing arterial occlusive disease like hypertension and elevated cholesterol. Five-year mortality after a diagnosis of CLTI is 60\% \citep{criqui1992mortality,van2018cardiovascular,baubeta2017amputation}. Poor systemic cardiovascular health increases the risk of heart attacks, stroke, and premature death, the severe arterial insufficiency in the arteries of the legs puts patients with CLTI patients at risk for non-healing foot ulcers, gangrene, or foot infections, all of which further elevate the risk of limb loss \citep{mills_society_2014, browder_analysis_2022, crowner_society_2022}. However, despite multiple treatment options, approximately half of patients with CLTI will experience limb loss or death within two years following surgical intervention \citep{farber_surgery_2022}. Our goal is to derive the optimal ITR for patients undergoing the two surgical treatments, endovascular therapy and open bypass surgery, by prioritizing OS while minimizing the risk of death, with major amputation before death considered as a competing risk.

Patients with more than one limb affected by disease were excluded from the study. Among those who experienced major amputation, we only included patients who had the amputation on the limb that underwent the index revascularization procedure. For instance, patients who had major amputation on the left limb but received open surgical revascularization on the right limb were excluded. Given that the decision to perform an open bypass is largely based on surgeon and patient preference, which may be influenced by the invasive nature of the procedure, we classified patients who initially underwent endovascular surgery (such as angioplasty or stent) but subsequently received an open bypass within a week as part of the open bypass treatment group. Our final analytic dataset had $N = 303$, with $77.9\%$ of patients in the endovascular treatment group and $22.1\%$ in the open bypass group. We analyzed 1-year follow-up data such that $\tau = 365$ days, and any observed events after $\tau$ were considered censored with $X_i = \tau$. The tailoring variables were age at treatment (in years), race, ethnicity, gender, severity of clinical presentation (limb-based Rutherford class), degree of ischemia, severity of foot wound anatomic location of arterial occlusive disease, previous revascularization history, and presences of clinical significant comorbidities including diabetes, hypertension, hyperlipidemia, renal disease, coronary artery disease, myocardial infarction, uncomplicated or complicated congestive heart failure,  cerebrovascular accident, obesity, smoking, chronic obstructive pulmonary disease, coagulopathy, anemia, dementia, venous insufficiency, venous thromboembolism (VTE), and unintentional weight loss. For more details, refer to the Supplementary Material.

\begin{figure}[H]
    \centering
    \includegraphics[width=0.9\textwidth]{Chapters/02_Proj1/03_figures/Figure_10CV_mean_rule1mean_surv_rule2gray_cr_tau365_2024-09-11.png}
    \caption{(Left) Estimated truncated mean survival time for preventing failure from all causes (major amputation before death or death) and (right) estimated truncated mean survival time for preventing mortality of the individualized treatment regime estimators. Each dot represents one cross-validation replicate. The black square dots are the average values of each method.}
	\label{fig:rda}
\end{figure}

Our results compared to the same methods in the simulation studies, the zero-order model, and the observed policy are presented in Figure \ref{fig:rda}. We illustrated the truncated mean failure time for both endpoints, where failure on the left is failure from any cause, and failure on the right is failure from death, the priority cause. AIPWE was unable to perform with the sample size less than $500$ for all CV iterations. Thus, we did not include all $K_{\text{CV}} = 10$ folds AIPWE in our results. From Figure \ref{fig:rda}, our method performs well compared to other existing precision medicine methods. We note that the zero-order model version of our method may outperform the precision medicine version of our method. However, this implies that our proposed two-phase method works well, but heterogeneity might not be enough in this dataset to distinguish between patients.  We have already show via simulations that heterogeneity is important to consider. The observed policy average also performed slightly better than our method average for mean truncated OS time. This suggested that the surgeon treatment recommendation is good for the patients, as expected from the experts in the field. Because our method also considers minimizing cumulative incidence of death, the average across folds for mean truncated death failure time is still higher than the observed policy. However, limitations in this analysis include an unbalanced number in the treatment groups and high censoring rate (70\%). 